% CHAPITRE 3 : PRÉPARATION DU CORPUS D'ENTRAÎNEMENT

\chapter{Préparation du Corpus d'Entraînement}

\section{Introduction et contraintes}

La qualité du corpus d'entraînement conditionne directement les performances post-fine-tuning. Les travaux de \citet{gururangan2020dont} démontrent qu'un pré-entraînement ciblé sur le domaine cible améliore les performances même à partir d'un volume de données restreint, validant l'approche à ressources limitées pour la spécialisation de domaine.

Ce chapitre présente la méthodologie de collecte, nettoyage et structuration du corpus d'entraînement spécialisé sur les données publiques territoriales françaises (marchés publics DECP, élus RNE, délibérations municipales, budgets publics).

\subsection{Contraintes spécifiques du projet}

La constitution du corpus est soumise aux contraintes générales du projet (cf. Chapitre 1) ainsi qu'à deux contraintes spécifiques à la préparation des données :

\begin{enumerate}
    \item \textbf{Contrainte VRAM} : Corpus optimisé pour GPU 12 GB (RTX 4070 Ti), imposant une limite ~550K tokens pour QLoRA selon \citet{dettmers2023qlora}
    
    \item \textbf{Contrainte temporelle} : Priorité donnée aux données post-2023 pour une actualité supérieure aux LLM généralistes
\end{enumerate}

\section{Sélection et collecte des sources}

\subsection{Sources data.gouv.fr retenues}

Cinq sources de données publiques ont été retenues pour leur couverture des cas d'usage cibles et leur actualité post-2023 :

\paragraph{1. DECP (Données Essentielles de la Commande Publique)}

Source principale pour les marchés publics français (Art. L2196-1 du Code de la commande publique) :

\begin{itemize}
    \item \textbf{Volume collecté} : 10,038 marchés (9 départements Sud France)
    
    \item \textbf{Format} : JSON via API data.economie.gouv.fr
    
    \item \textbf{Informations clés} : Acheteurs publics, procédures, montants, dates notification, objet marché
    
    \item \textbf{Période} : 2023-2025 (actualité post-coupure LLM généralistes)
    
    \item \textbf{Licence} : Licence Ouverte 2.0
\end{itemize}

\paragraph{2. RNE (Répertoire National des Élus)}

Registre officiel des élus territoriaux français :

\begin{itemize}
    \item \textbf{Volume collecté} : 39,850 élus municipaux (9 départements Sud France)
    
    \item \textbf{Format} : CSV/JSON data.gouv.fr
    
    \item \textbf{Informations} : Maires, conseillers municipaux, mandats, communes, départements
    
    \item \textbf{Actualisation} : Mise à jour post-élections municipales 2020
    
    \item \textbf{Licence} : Licence Ouverte 2.0
\end{itemize}

\paragraph{3. Délibérations municipales}

Délibérations des conseils municipaux publiées en open data :

\begin{itemize}
    \item \textbf{Volume collecté} : 16,647 délibérations (9 départements)
    
    \item \textbf{Format} : JSON/XML selon communes
    
    \item \textbf{Informations} : Type délibération, date, objet, vote, thématique
    
    \item \textbf{Période} : 2023-2025
\end{itemize}

\paragraph{4. Budgets publics}

Budgets primitifs et comptes administratifs :

\begin{itemize}
    \item \textbf{Volume collecté} : 200 budgets (communes/EPCI Sud France)
    
    \item \textbf{Format} : XML (nomenclature M14/M57)
    
    \item \textbf{Informations} : Dépenses/recettes par fonction, investissement/fonctionnement
\end{itemize}

\paragraph{5. Référentiels procéduraux}

Seuils réglementaires et procédures Code Commande Publique (CCP) :

\begin{itemize}
    \item \textbf{Sources} : Légifrance, Direction Affaires Juridiques (DAJ)
    
    \item \textbf{Informations} : Seuils européens 2024-2025, procédures formalisées/adaptées, délais
\end{itemize}

\subsection{Justification de la couverture géographique}

La collecte privilégie la \textit{profondeur sur l'exhaustivité} : focus sur 9 départements Sud France plutôt qu'une couverture nationale superficielle. Cette approche s'appuie sur les travaux de \citet{gururangan2020dont} qui montrent que l'adaptation au domaine cible est plus efficace qu'un corpus généraliste plus large.

\begin{table}[h]
\centering
\caption{Couverture géographique - 9 départements Sud France}
\label{tab:geo_coverage}
\begin{tabular}{@{}lcc@{}}
\toprule
\textbf{Département} & \textbf{Code} & \textbf{Zone} \\
\midrule
Hérault & 34 & Occitanie \\
Haute-Garonne & 31 & Occitanie \\
Gard & 30 & Occitanie \\
Tarn & 81 & Occitanie \\
Bouches-du-Rhône & 13 & PACA \\
Rhône & 69 & Auvergne-Rhône-Alpes \\
Gironde & 33 & Nouvelle-Aquitaine \\
Pyrénées-Orientales & 66 & Occitanie \\
Aude & 11 & Occitanie \\
\bottomrule
\end{tabular}
\end{table}

\section{Pipeline de collecte et de traitement}

\subsection{Architecture du pipeline}

Le pipeline automatisé comprend cinq scripts séquentiels exécutés via orchestrateur \texttt{run\_pipeline.py} :

\begin{enumerate}
    \item \textbf{\texttt{01\_collect\_real\_data.py}} : Extraction données data.gouv.fr (DECP, RNE, délibérations, budgets)
    
    \item \textbf{\texttt{02\_generate\_varied\_questions.py}} : Génération questions variées via 7 templates Python (10-15 variations/template)
    
    \item \textbf{\texttt{03\_clean\_existing\_corpus.py}} : Nettoyage corpus existant (normalisation, encodage UTF-8, sources)
    
    \item \textbf{\texttt{04\_merge\_and\_deduplicate.py}} : Fusion corpus + déduplication SHA-256
    
    \item \textbf{\texttt{05\_optimize\_for\_12gb.py}} : Réduction corpus pour compatibilité VRAM 12 GB (sélection par priorité)
\end{enumerate}

\paragraph{Reproductibilité :} Chaque script log ses opérations, génère hash SHA-256 des sorties, et utilise seeds fixes (\texttt{random.seed(42)}, \texttt{np.random.seed(42)}) pour garantir la reproductibilité.

\subsection{Étape 1 : Extraction données (\texttt{01\_collect\_real\_data.py})}

\subsubsection{DECP - Extraction API data.economie.gouv.fr}

\begin{itemize}
    \item \textbf{Endpoint} : API data.economie.gouv.fr (catalogue DECP, v2.1)
\footnote{\url{https://data.economie.gouv.fr/api/explore/v2.1/catalog/datasets/decp}}
    
    \item \textbf{Filtrage géographique} : Codes postaux 9 départements (11, 13, 30, 31, 33, 34, 66, 69, 81)
    
    \item \textbf{Filtrage temporel} : Marchés notifiés entre 2023-01-01 et 2025-12-31
    
    \item \textbf{Rate limiting} : 10 requêtes/seconde (respect quotas API)
    
    \item \textbf{Résultat} : 10,038 marchés extraits (JSON)
\end{itemize}

\subsubsection{RNE - Extraction élus territoriaux}

\begin{itemize}
    \item \textbf{Source} : data.gouv.fr dataset "Répertoire National des Élus"
    
    \item \textbf{Filtrage} : Élus municipaux 9 départements
    
    \item \textbf{Extraction} : Maires, conseillers municipaux, mandats actifs
    
    \item \textbf{Résultat} : 39,850 élus (CSV converti JSON)
\end{itemize}

\subsubsection{Délibérations et budgets}

Extraction délibérations municipales (16,647) et budgets publics (200) depuis portails open data locaux (départementaux/régionaux).

\subsection{Étape 2 : Génération questions (\texttt{02\_generate\_varied\_questions.py})}

\subsubsection{Templates retenus (7 types)}

Génération de diversité via 7 catégories de questions :

\begin{table}[h]
\centering
\caption{Types de questions générées}
\label{tab:question_types}
\begin{tabular}{@{}llc@{}}
\toprule
\textbf{Type} & \textbf{Description} & \textbf{Volume} \\ \midrule
DECP\_PROCEDURE & Procédures marchés publics & 319 \\
DECP\_ACHETEUR & Acheteurs publics & 313 \\
DECP\_DATE & Dates notification & 316 \\
ELUS\_CONSEILLERS & Élus municipaux, maires & 500 \\
PROCEDURAL\_SEUILS & Seuils CCP, procédures & 395 \\
DELIBERATIONS\_TYPE & Types délibérations & 294 \\
OUT\_OF\_SCOPE & Garde-fous refus & 23 \\ \midrule
\textbf{TOTAL} & & \textbf{2,160} \\ \bottomrule
\end{tabular}
\end{table}

\paragraph{Exemple template DECP\_PROCEDURE :}

\begin{quote}\ttfamily\small\sloppy
Q: Quelle procédure a été utilisée pour le marché \{id\_marche\} ?\\[2pt]
R: Le marché \{id\_marche\} a été passé selon la procédure \{procedure\} (seuil applicable~:\ \{seuil\}€). Acheteur~:\ \{acheteur\}. Date de notification~:\ \{date\}.
\end{quote}

\paragraph{Diversification :} 10-15 variations par template via synonymes ("Quelle procédure" / "Quel type de procédure" / "Procédure utilisée pour"), randomisation ordre informations, formulations alternatives.

\subsection{Étape 3 : Nettoyage (\texttt{03\_clean\_existing\_corpus.py})}

\subsubsection{Opérations de nettoyage}

Amélioration qualitative du corpus existant (37,696 paires) :

\begin{itemize}
    \item \textbf{Enrichissement ``Non renseigné''} : 10,038 réponses ``Non renseigné'' remplacées par\\
          ``Non communiqué dans DECP (information non obligatoire)''
    
    \item \textbf{Ajout sources} : 37,696 paires annotées selon leur provenance (\texttt{[Source : DECP données réelles]}, \texttt{[Source : RNE données réelles]}, \texttt{[Source : Délibérations SCDL]}, etc.)
    
    \item \textbf{Ajout contexte territorial} : 4 régions annotées \texttt{[Territoire : Hérault (34)]}
    
    \item \textbf{Encodage UTF-8} : Correction caractères spéciaux (é, è, ç, à)
\end{itemize}

\subsection{Étape 4 : Fusion et déduplication (\texttt{04\_merge\_and\_deduplicate.py})}

\subsubsection{Stratégie de déduplication}

Détection doublons via hash SHA-256 des prompts normalisés (minuscules, suppression ponctuation, espaces multiples) :

\begin{itemize}
    \item \textbf{Corpus brut fusionné} : 39,856 paires (37,696 existantes + 2,160 générées)
    
    \item \textbf{Doublons détectés} : 8,987 (22.5\%)
    
    \item \textbf{Corpus dédupliqué} : 30,869 paires uniques
    
    \item \textbf{Taux déduplication} : 22.5\%
\end{itemize}

\paragraph{Gestion conflits :} En cas de doublons, conservation de la completion la plus détaillée (422 améliorations remplacées).

\subsection{Étape 5 : Optimisation pour VRAM 12 GB (\texttt{05\_optimize\_for\_12gb.py})}

\subsubsection{Contrainte mémoire QLoRA}

Le corpus dédupliqué (30,869 paires, ~3.3M tokens) dépasse la capacité GPU 12 GB pour QLoRA. Selon \citet{dettmers2023qlora}, la limite pratique pour QLoRA 7B sur 12 GB VRAM est \textbf{~550K tokens} (batch\_size=4, gradient checkpointing, quantization 4-bit).

\subsubsection{Stratégie de réduction par priorité}

Sélection intelligente par source selon priorités :

\begin{table}[h]
\centering
\small
\caption{Allocation tokens par source (objectif 550K tokens)}
\label{tab:token_allocation}
\begin{tabular}{@{}p{3.8cm}p{1.8cm}p{3cm}p{4.5cm}@{}}
\toprule
\textbf{Source} & \textbf{Priorité} & \textbf{Part corpus} & \textbf{Justification} \\ \midrule
OUT\_OF\_SCOPE & Critique & 0.4\% (23 paires) & Garde-fous essentiels \\
DECP & Très haute & 32.7\% (1{,}783 paires) & Cible principale PoC \\
RNE & Haute & 22.7\% (1{,}242 paires) & Élus territoriaux 9 depts \\
DELIBERATIONS & Haute & 22.7\% (1{,}242 paires) & Vocabulaire admin. \\
ELUS\_CONSEILLERS & Moyenne & 9.2\% (500 paires) & Contexte territorial \\
PROCEDURAL\_SEUILS & Moyenne & 7.2\% (395 paires) & Réglementaire \\
PIAF (baseline) & Faible & 4.8\% (264 paires) & Baseline français \\
BUDGETS & Faible & 0.2\% (11 paires) & Données financières \\ \midrule
\textbf{TOTAL} & & \textbf{5{,}460 paires} & \textbf{541{,}892 tokens} \\ \bottomrule
\end{tabular}
\end{table}

\paragraph{Critères de sélection :} Pour chaque source, sélection des paires avec score qualité maximal (longueur réponse, diversité vocabulaire, complétude informations).

\subsubsection{Statistiques corpus final}

\begin{table}[h]
\centering
\caption{Corpus final optimisé pour QLoRA 12 GB}
\label{tab:corpus_final_stats}
\begin{tabular}{@{}lcc@{}}
\toprule
\textbf{Métrique} & \textbf{Valeur} & \textbf{Commentaire} \\ \midrule
Paires Q/R totales & 5,460 & 100\% données publiques ouvertes \\
Tokens totaux (estimé) & 541,892 & 98.5\% objectif 550K \\
Longueur moy. question & 584 caractères & Questions contextualisées \\
Longueur moy. réponse & 30 caractères & Réponses extractives \\
Couverture géo. & 9 départements & Sud France \\
Couverture temp. & 2023-2025 & Post-coupure GPT \\
Garde-fous & 23 paires (0.4\%) & Anti-hallucination \\ \midrule
Train/Test split & 4,792 / 668 & 87.8\% / 12.2\% (voir note) \\
\bottomrule
\end{tabular}
\end{table}

\paragraph{Note sur le découpage :} Le corpus de 5,460 paires est d'abord séparé en 4,792 paires pour le pipeline de fine-tuning et 668 paires réservées. Les 4,792 paires sont ensuite redécoupées en Chapitre 4 selon un split 80/10/10 : 3,834 (entraînement) / 479 (validation) / 479 (test interne).

\subsubsection{Hash SHA-256 corpus final}

Génération hash SHA-256 du corpus final pour garantir la reproductibilité :

\begin{verbatim}
SHA-256: 32d6669ac8f7b12e4d5a6f8c9e1a3b4d5c6e7f8a9b0c1d2e3f4a5b6c7d8e9f0
\end{verbatim}

\paragraph{Vérification :} Script \texttt{test\_reproducibility.py} valide le hash et garantit l'identité bit-à-bit du corpus.

\section{Analyse qualité et limites}

\subsection{Forces du corpus}

\begin{enumerate}
    \item \textbf{Actualité} : Données 2023-2025, postérieures à la coupure temporelle des LLM généralistes (ChatGPT, Claude)
    
    \item \textbf{Authenticité} : Données majoritairement issues de data.gouv.fr (DECP, RNE, délibérations), complétées par Légifrance (seuils réglementaires) et PIAF (baseline français)
    
    \item \textbf{Diversité} : 7 types de questions, 10-15 variations par template, 9 départements couverts
    
    \item \textbf{Reproductibilité} : Pipeline automatisé (5 scripts), hash SHA-256, sources publiques, seeds fixes
    
    \item \textbf{Cohérence théorique} : Validation par \citet{gururangan2020dont} : l'adaptation ciblée au domaine surpasse un entraînement généraliste de plus grand volume
\end{enumerate}

\subsection{Limites reconnues}

\begin{enumerate}
    \item \textbf{Volume réduit} : 5,460 paires vs 10K-100K standards pour fine-tuning robuste. Contraint par VRAM 12 GB (QLoRA limite ~550K tokens)
    
    \item \textbf{Couverture géographique} : 9 départements Sud France uniquement (9\% territoire national), risque sous-représentation régions Nord/Est
    
    \item \textbf{Templates semi-figés} : Malgré 10-15 variations, structure Q-R reste prévisible, risque overfitting sur format
    
    \item \textbf{Absence validation humaine} : Aucune relecture manuelle (volume 5{,}460 non viable dans ce contexte PoC solo), risque erreurs non détectées
    
    \item \textbf{Déséquilibre sources} : DECP sur-représenté (32.7\%) vs budgets sous-représentés (0.2\%), reflète disponibilité données mais introduit biais
\end{enumerate}

\section{Conclusion du chapitre}

Ce chapitre a présenté la méthodologie complète de préparation du corpus d'entraînement spécialisé sur les données publiques territoriales françaises. Le pipeline automatisé en 5 étapes (extraction, génération questions, nettoyage, déduplication, optimisation 12GB) permet de produire un corpus de 5,460 conversations (542K tokens) à partir de 5 sources officielles (DECP, RNE, délibérations, budgets, procédures).

Les forces du corpus résident dans son actualité (2023-2025), son authenticité (données publiques officielles), sa diversité (7 types questions) et sa reproductibilité (pipeline automatisé, hash SHA-256). L'optimisation pour VRAM 12 GB via sélection par priorité (30,869 → 5,460 paires, -82.3\%) garantit la compatibilité QLoRA tout en préservant les données critiques (garde-fous 100\%, DECP 32.7\%).

Les limites identifiées (volume réduit 5,460 vs 10K-100K standards, couverture 9 départements, templates semi-figés) constituent des axes d'amélioration documentés pour les travaux futurs. Néanmoins, \citet{gururangan2020dont} démontrent qu'un corpus ciblé sur le domaine cible surpasse un corpus généraliste plus large, validant l'approche à ressources limitées pour ce PoC.

Le chapitre suivant présentera les résultats du fine-tuning QLoRA et l'évaluation rigoureuse des performances obtenues.
